\section{NVIDIA Jetson}

\subsection{Jetson Nano}

Jetson Nano đại diện cho nhóm thiết bị Edge AI có tài nguyên hạn chế. Thiết bị
phù hợp để đánh giá mức tối thiểu mà model có thể vận hành, đồng thời làm rõ ảnh
hưởng của bộ nhớ, power mode và thermal throttling. Phiên bản JetPack, CUDA,
TensorRT, dung lượng RAM và chế độ công suất phải được ghi theo đúng thiết bị
thực nghiệm, vì các biến thể phần cứng có thể khác nhau.

\subsection{Jetson Orin Nano}

Jetson Orin Nano đại diện cho thế hệ GPU nhúng mới hơn, có kiến trúc và năng lực
tăng tốc khác Jetson Nano. Thiết bị được dùng để đánh giá khả năng mở rộng hiệu
năng của cùng pipeline model, đồng thời kiểm tra tính di động của Model Registry,
Deployment và Edge Runtime giữa hai lớp node.

\subsection{So sánh tài nguyên phần cứng}

\begin{table}[H]
\centering
\caption{Thông tin phần cứng cần ghi nhận cho hai thiết bị Jetson}
\label{tab:jetson-hardware}
\begin{tabularx}{\textwidth}{L{4cm}YY}
\toprule
\textbf{Thành phần} & \textbf{Jetson Nano} & \textbf{Jetson Orin Nano}\\
\midrule
Kiến trúc GPU & \cbs{thông tin thiết bị} & \cbs{thông tin thiết bị}\\
Số nhân GPU/Tensor Core & \cbs{thông tin thiết bị} & \cbs{thông tin thiết bị}\\
RAM & \cbs{dung lượng} & \cbs{dung lượng}\\
JetPack & \cbs{phiên bản} & \cbs{phiên bản}\\
CUDA & \cbs{phiên bản} & \cbs{phiên bản}\\
TensorRT & \cbs{phiên bản} & \cbs{phiên bản}\\
Power mode & \cbs{chế độ đo} & \cbs{chế độ đo}\\
Tản nhiệt & \cbs{cấu hình} & \cbs{cấu hình}\\
\bottomrule
\end{tabularx}
\end{table}

Bảng~\ref{tab:jetson-hardware} phải được điền từ lệnh hệ thống và nhãn thiết bị
thực tế, không suy ra từ thông số quảng cáo của một phiên bản khác.

\subsection{CUDA và GPU Acceleration}

CUDA cung cấp mô hình lập trình và runtime để thực thi kernel trên GPU NVIDIA.
Trong pipeline, TensorRT sử dụng CUDA để lập lịch kernel và truyền dữ liệu. Tối
ưu end-to-end cần hạn chế sao chép CPU--GPU, tái sử dụng bộ đệm, dùng stream
phù hợp và đồng bộ chỉ tại những điểm cần thiết.

GPU acceleration không tự động bảo đảm latency thấp. Kích thước input, batch,
memory allocation, camera decode và hậu xử lý trên CPU có thể trở thành nút thắt.
Vì vậy, đồ án đo từng giai đoạn của pipeline và tổng thời gian đầu cuối.